\section{本次运行的数值记录}
下表由 JSON 记录生成。CPU 环境为 Intel Core Ultra 7 155H，NumPy 1.26.4、SciPy 1.15.3，双精度、单线程 BLAS。时间仅代表这些实现和小规模工作负载。
\begin{center}\small
\begin{tabular}{lrrrrr}\toprule
案例 & 网格/单元数 & $p$ & 步数 & $L^2$ 误差 & 推进时间/s\\\midrule
Yee & 12 & -- & 14 & 2.076e-02 & 1.424e-03\\
Yee & 24 & -- & 27 & 5.094e-03 & 8.223e-03\\
DG & 8 & 1 & 44 & 2.523e-02 & 2.015e-03\\
DG & 16 & 1 & 87 & 6.550e-03 & 4.003e-03\\
DG & 8 & 2 & 120 & 1.665e-03 & 5.416e-03\\
DG & 16 & 2 & 240 & 2.091e-04 & 1.100e-02\\
DG & 8 & 3 & 236 & 7.970e-05 & 1.056e-02\\
DG & 16 & 3 & 471 & 5.295e-06 & 2.169e-02\\
FVM & 32 & 0 & 12 & 5.534e-02 & 6.787e-05\\
FVM & 64 & 0 & 24 & 2.786e-02 & 1.272e-04\\
\bottomrule\end{tabular}\end{center}
Yee 的终止时间为 $T=0.5$，DG 与 FVM 为 $T=0.15$。DG 与 FVM 均以重构场对解析场的积分误差比较；另行保存 FVM 单元平均误差。不同方程的误差和时间分别解读。
\section{GPU 与 MPI 执行记录}
GPU 采用常系数融合更新（30 FLOPs/胞元整步）；设备推进与包含传输的时间分别报告。两种网格均固定 30 步，终止时间不同，故不由其误差比估计收敛阶。远端 GPU 为共享负载下的短时验证。
\begin{center}\small\begin{tabular}{lrrrrr}\toprule
设备 & $n$ & $T$ & 核推进/s & 含传输/s & CPU 最大差\\\midrule
本机 RTX 500 Ada & 32 & 0.3750 & 5.591e-04 & 1.427e-03 & 3.37e-16\\
本机 RTX 500 Ada & 64 & 0.1875 & 3.622e-03 & 1.015e-02 & 4.44e-16\\
GK RTX 4090 & 32 & 0.3750 & 1.745e-04 & 8.856e-04 & 3.37e-16\\
GK RTX 4090 & 64 & 0.1875 & 4.792e-04 & 5.303e-03 & 4.44e-16\\
\bottomrule\end{tabular}\end{center}
MPI 实验均在单一主机内运行，包含 halo 交换，不能用作跨节点扩展数据。
\begin{center}\small\begin{tabular}{llrrrr}\toprule
位置 & 案例 & 网格/单元数 & 进程 & $L^2$ 误差 & 最慢进程时间/s\\\midrule
本机 & Yee & 12 & 2 & 2.076e-02 & 2.516e-04\\
本机 & Yee & 24 & 2 & 5.094e-03 & 3.299e-03\\
本机 & DG($p=3$) & 8 & 2 & 7.970e-05 & 4.386e-04\\
本机 & DG($p=3$) & 16 & 2 & 5.295e-06 & 7.949e-04\\
GK & Yee & 12 & 2 & 2.076e-02 & 1.304e-04\\
GK & Yee & 24 & 2 & 5.094e-03 & 1.768e-03\\
GK & DG($p=3$) & 8 & 2 & 7.970e-05 & 2.962e-04\\
GK & DG($p=3$) & 16 & 2 & 5.295e-06 & 6.748e-04\\
\bottomrule\end{tabular}\end{center}
一进程和两进程的完整记录、守恒与约束量见相应硬件 JSON。本次记录未使用 DRAM 计数器，也未标定硬件 Roofline；语义流量与理想复用估算在正文中单列。
